\section{Trabalhos Relacionados}
\label{sec:related_work}

O estudo de \cite{Bruno2024} propõe uma metodologia para avaliar o nível de acessibilidade a serviços essenciais em cidades ao redor do mundo a partir do conceito de cidade de 15 minutos. A pesquisa mostra que mesmo em cidades com infraestrutura razoável, a distribuição desigual de serviços e a variação na densidade populacional geram disparidades significativas no acesso. Ao quantificar essas diferenças e propor cenários de redistribuição, os autores oferecem contribuições relevantes para análises de acessibilidade urbana. Esses achados se conectam diretamente com os objetivos deste trabalho, que busca identificar desigualdades espaciais no acesso a serviços em Curitiba.

\cite{Nicoletti2022} apresentam uma análise baseada em dados abertos para investigar como diferentes perfis demográficos acessam a infraestrutura urbana. O estudo, feito em dez grandes cidades do mundo, mostra que comunidades mais vulneráveis têm menos acesso a serviços essenciais. A metodologia cruza dados de infraestrutura e dados socioeconômicos, permitindo identificar de forma clara as desigualdades espaciais. Este trabalho é relevante pois oferece uma base para compreender como condições socioeconômicas distintas impactam a acessibilidade espacial. Também destaca que a integração de análises espaciais e socioeconômicas podem revelar dinâmicas internas de desigualdade em centros urbanos.

\cite{Adorno2025} combinam técnicas de clusterização espacial e regressão geograficamente ponderada (GWR) para analisar desigualdades na acessibilidade a áreas verdes urbanas em Goiânia. O estudo identifica agrupamentos de baixa acessibilidade e correlaciona-os com variáveis socioeconômicas, evidenciando como características demográficas e territoriais impactam o acesso a espaços públicos. A abordagem metodológica proposta é altamente relevante para este trabalho, que também busca investigar a relação entre desigualdades socioeconômicas e acessibilidade urbana por meio de técnicas espaciais.

Ainda, \cite{Araujo2025} analisam como gênero e renda afetam a acessibilidade em Curitiba por meio de simulações individuais baseadas em agentes. O estudo mostra que mulheres de baixa renda têm mais dificuldade de acesso a oportunidades, mesmo em regiões centrais. A abordagem revela desigualdades que não aparecem em análises convencionais, o que reforça a importância de considerar múltiplos recortes sociais neste trabalho. Considera-se a possibilidade de incorporar, futuramente, abordagens similares de simulação baseada em agentes para aprofundar a compreensão das dinâmicas locais de acessibilidade.

Embora os estudos revisados compartilhem o interesse por desigualdades na acessibilidade urbana, este trabalho foca especificamente em acessibilidade temporal a pé e na integração de métodos globais (regressão) e locais (clusterização espacial). Essa abordagem permite capturar tanto padrões gerais da relação entre fatores socioeconômicos e acessibilidade quanto heterogeneidades intramunicipais, oferecendo um diagnóstico espacial com potencial para orientar políticas urbanas precisas, superando limitações de análises puramente teóricas ou baseadas em simulações.